\section{A common prime list for growing families of different moduli}\label{sec:mixed}
We now allow a separate modulus for every block and use every retained prime. A block of modulus one plays the role of a reserve. The open question on growing families following \texttt{dir:thm:minority} asks whether relative counts on occupied fibres survive below majority retention. They survive up to the full order of growth allowed by the prime cutoff and the quotas.

\begin{theorem}[Independent moduli after positive-density retention]\label{th:mixed}
Fix \(0<a\le1\) and \(A>0\), and put \(D=\lfloor1/a\rfloor\). There are positive constants \(B,K_1,K_2,c\), depending only on \(a,A\), with the following property under Dirichlet GRH. Write \(L=\log Q\). For every sufficiently large \(Q\) and every integer \(2\le R\le cL/\log L\), there is a single list \(\mathcal S\) of between \(K_1 R L\) and \(2K_1 R L+2\) distinct primes in \([x,2x]\), where \(x\asymp_a L^2\) and \(x\le K_2 L^2\).

Choose \(2\le r\le R\) and \(1\le q_j\le Q\) independently. Set \(M=\operatorname{lcm}_j q_j\), \(g_0=\gcd_j q_j\), and \(U=\{p\in\mathcal S:p\nmid M\}\). Retain any \(T\subseteq U\) of size \(m\ge a|U|\). For a uniform partition of \(T\) into \(r\) labelled blocks of sizes \(n_j\) with \(h=\min_j n_j\ge BL\), let \(\mu\) be the law of their product vector in \(V=\prod_j(\mathbb Z/q_j\mathbb Z)^\times\). Balanced admissible quotas exist for each such \(T\).

There is a subgroup \(E\subseteq\widehat V\), independent of the quotas, containing the diagonal character group \(\Delta\simeq\widehat{(\mathbb Z/g_0\mathbb Z)^\times}\), and
\[
 [E:\Delta]\le D^{r-1}.
\]
For evaluation \(\Pi:V\to\widehat E\), let \(\nu=\Pi_*\mu\). The support of \(\mu\) is precisely \(\Pi^{-1}(\operatorname{supp}\nu)\). On it,
\[
   \mu(v)=\frac{|E|}{|V|}\nu(\Pi(v))(1+\theta_v),
       \qquad |\theta_v|\le Q^{-A}.
\]
The projected support has at most \([E:\Delta]\) points. For \(a>1/2\) the sole constraint is that the product of the target residues agree with \(\prod_{p\in T}p\) modulo \(g_0\).
\end{theorem}

\begin{proof}
\emph{A list that can test short products of characters.}
Choose \(\tau,\eta>0\) so that \(1-D^2\tau>0\) and
\((D+1)a(1-D^2\tau)^2>1+D\eta\), and set \(\delta=\eta/4\), reducing \(\eta\) if needed to be less than one. We use the source's common-list version of the prime-pool estimate \eqref{eq:pool} (the one-set lemma in \cite{factor}, labelled \texttt{dir:lem:oneset} in the enlarged article). With parameter \(Q_*=Q^{2D}\) it supplies, for every integer length parameter
\[
  C_\delta\log(2Q_*)\le N_*\le c_\delta\log^2(2Q_*)/\log\log(3Q_*),
\]
a single list of between \(N_*\) and \(2N_*\) primes in \([x,2x]\), \(x=K_\delta\log^2(2Q_*)\), on which every nonprincipal character of modulus at most \(Q_*^2=Q^{4D}\) has mean of modulus at most \(\delta\). The source derives the dyadic estimate by the smoothed GRH explicit formula and selects the list by independent sampling.

Take \(B\ge \max\{1,(2/\tau)(A+36D^2+13)\}\). Next choose \(K_1\) large enough for the lower bound on \(N_*=\lceil K_1RL\rceil\) and \(aK_1/2\ge 2(B+1)\). Finally choose \(0<c\le1\) small enough for the upper bound throughout the range of \(R\). All choices are independent of \(q_j,T,r\). At most \(rL/\log x\) list primes divide \(M\), as follows by taking their product. This is \(o(|\mathcal S|)\) uniformly, so every indicated character has mean on \(U\) of modulus at most \(\eta\) for large \(Q\), and \(m\ge a K_1 R L/2\). This proves the balanced-quota assertion. Also, since \(r\le L/\log L\) and \(m\le 2K_1 R L+2\), for large \(Q\) one has
\begin{equation}\label{eq:mixedsize}
            r\log(m+1)\le 3L .
\end{equation}
The coefficient three here is independent of the fixed choice of \(K_1\).

\emph{The subgroup seen by the available characters.}
Inflate characters of each block to the unit group modulo \(M\), writing the inflated character as \(\widetilde\psi_j\). Generate \(J\) by all ratios \(\widetilde\psi_i\overline{\widetilde\psi_j}\) of block characters whose bias on \(T\) exceeds \(1-\tau\), allowing all choices of the two characters. To control this subgroup, take two different words of length at most \(D\) in the generators and inverses, viewed as different characters. Their quotient involves at most \(4D\) block characters counting repetitions. It is induced by a nonprincipal character modulo the least common multiple of the participating moduli, an integer at most \(Q^{4D}\). On \(U\) the induced and inflated values coincide. Thus the corresponding distinct word characters have Gram off-diagonal entries bounded by \(\eta\), even though \(M\) can be larger than any fixed power of \(Q\). Each word has bias at least \(1-D^2\tau\) on \(T\) by the telescoping estimate used after \eqref{eq:pool}. The same Gram-matrix and word-ball argument with \(|T|/|U|\ge a\) now gives \(|J|=d\le D\).

Define \(E\) to consist of tuples \((\psi_1,\dots,\psi_r)\) with all inflated characters in a single coset of \(J\). This is a subgroup. Taking their ratios with \(\widetilde\psi_1\) gives a homomorphism \(E\to J^{r-1}\) whose kernel consists of equal inflated characters. A common character descends to the gcd of the moduli: its unique primitive conductor divides every \(q_j\). Conversely every character modulo \(g_0\) inflates to such a tuple. This identifies the kernel with \(\Delta\) and proves the index bound, including unit moduli.

\emph{Uniform Fourier decay.}
For a character tuple outside \(E\), let \(t\ge1\) be the number of labels outside a coset of largest combined quota. Distinct cosets have bias gap at least \(\tau\) on all ratios of their corresponding labels, by construction of \(J\). The pair estimate retaining all cross-coset pairs in Theorem~\ref{th:growing} bounds that coefficient by \((m+1)^r e^{-\tau t h/2}\). There are at most
\[
 rQ\binom rt d^r Q^t
\]
such tuples for each \(t\): specify one character in the dominant coset from some block, the outside labels, and their characters. Each inside label has at most \(d\) possibilities since inflation from its modulus is injective. Since \(rQ e^{-\tau h/2}\le1/2\) for large \(Q\), summing as in \eqref{eq:growerror} gives
\[
 \sum_{\boldsymbol\psi\notin E}|\widehat\mu(\boldsymbol\psi)|
 \le (m+1)^r rQ d^r\big[(1+Qe^{-\tau h/2})^r-1\big]
 \le 2Q^{10-\tau B/2}.
\]
We used \eqref{eq:mixedsize}, \(d\le m+1\), and \(r\le Q\). Fourier inversion over \(E\) contributes exactly \(|E|\nu(\Pi(v))\) to \(|V|\mu(v)\).

\emph{Mass of an occupied projection.}
Color each prime by its value in \(\widehat J\) under evaluation. For a tuple in \(E\), put \(\delta_j=\widetilde\psi_j/\widetilde\psi_1\in J\). If the blocks have integer products \(P_j\), the value of the tuple is
\[
 \widetilde\psi_1\Big(\prod_{p\in T}p\Big)\prod_{j=1}^r\delta_j(P_j).
\]
Consequently preserving the \(r\) color products of any witness preserves its complete \(E\)-projection. The integer-table repair in Theorem~\ref{th:growing}, with \(r\) rows and \(d\) colors, realizes the witness's color products from every allocation using at most \(6D^2r\) edited prime labels. Hence each occupied projection has mass at least
\[
    (1+mr)^{-6D^2r}\ge Q^{-36D^2}
\]
by \eqref{eq:mixedsize}. The Fourier error is at most \(Q^{-A}\) times every occupied main term for large \(Q\) by our choice of \(B\). Empty projections cannot occur. All occupied projections restrict to the same value on \(\Delta\), determined by the total product. This gives the bound on their number and the assertion for \(D=1\).
\end{proof}

The order of the growing range is sharp for the stated size of primes and balanced-quota guarantee. Even with no deletions, \(r\) blocks of size \(B\log Q\) use at least \(rB\log Q\) distinct primes, whereas \([1,2K_2(\log Q)^2]\) contains only \(O_{K_2}((\log Q)^2/\log\log Q)\) primes. Thus the loss allowed by the quotient introduces no further restriction in order of magnitude on the number of independently prescribed moduli.
